%% ============================================================
%% preamble/commands.tex — semantic inline macros
%%
%% Purpose      : implements the authoring contract's inline
%%                vocabulary (docs/architecture/authoring-contract.md).
%%                These are SEMANTIC hooks: print styling lives here,
%%                and the EPUB converter maps each macro to a class
%%                (\term -> <i class="term">, etc.). Chapters may use
%%                these macros and nothing else.
%% Dependencies : etoolbox; code.tex (which owns \code when the code
%%                module is on — hence this file loads after it).
%% Provenance   : the-last-book's PARSE CONTRACT macro table
%%                (docs/research/the-last-book.md).
%% ============================================================

% \term — first use of a technical term (italic; converter: i.term).
\newcommand{\term}[1]{\emph{#1}}

% \work — titles of books/films/periodicals (italic; converter: i.work).
% Distinct from \emph so the EPUB can class them differently and an
% index pass can find every cited work.
\newcommand{\work}[1]{\emph{#1}}

% \person — personal names. Print styling is deliberately plain by
% default; books that want Bringhurst-style small-caps introductions
% swap the definition here (never in chapters). Converter: span.person.
\newcommand{\person}[1]{#1}

% \keyterm — a defined term at its definition site (bold; converter:
% dfn.keyterm). Indexing is deliberately separate: not every locally defined
% phrase is a durable lookup term, and an automatic index is a concordance.
% Use \indexentry at definitions and substantive later discussions.
\newcommand{\keyterm}[1]{\textbf{#1}}
\newcommand{\indexentry}[1]{\index{#1}}

% \foreignphrase — non-English phrases (italic; converter: i.foreign,
% which also earns a lang attribute during EPUB conversion).
\newcommand{\foreignphrase}[1]{\emph{#1}}

% \figalt — accessibility alternative text for the enclosing figure.
% No-op in print (alt text is an EPUB concern); the converter uses it
% as the image's alt attribute. \figalt{} (empty) marks the image
% decorative (alt="" role="presentation"). Without \figalt the
% converter falls back to the caption text; a figure with neither is
% an epub-check violation. Keep it under ~140 characters (KDP).
\newcommand{\figalt}[1]{}

% \code — inline code. The code module (preamble/code.tex) defines the
% robust \lstinline-backed version; this fallback keeps the macro
% available (simple \texttt) in books with modules.code: false so
% prose mentioning \code{make pdf} still compiles.
\ifBookModuleCode\else
  \newcommand{\code}[1]{\texttt{#1}}
\fi

% ----------------------------------------------------------------------
% Block-quote support: \attribution closes a quotation with a
% right-aligned em-dash credit line, e.g.
%   \begin{quotation} ... \attribution{Robert Bringhurst} \end{quotation}
% The \nopagebreak keeps a credit from stranding on the next page.
% ----------------------------------------------------------------------
\newcommand{\attribution}[1]{%
  \par\nopagebreak\smallskip
  {\raggedleft\normalfont\color{bk-attribution}\textemdash\,#1\par}}

% Block quotes sit a touch smaller than body copy (trade convention);
% the size change also signals "quotation" in grayscale POD where no
% color cue exists.
\AtBeginEnvironment{quotation}{\small}
